% Cardiac MRI LGE Patterns - 心臟磁振造影 Late Gadolinium Enhancement 判讀教學
% 作者: 謝慕揚, MD, PhD, FESC
% 圖片來源: www.radiologyillustration.com
\documentclass[12pt,a4paper]{article}
% Drake_preamble.tex - LaTeX Preamble Template
% 作者: 謝慕揚, MD, PhD, FESC
% 
% 使用說明:
% 1. 表格請使用 longtable，不要有垂直分隔線 |
% 2. 表格內換行使用 \newline，不要使用 \\
% 3. 藥物名稱、檢驗、檢查請維持英文
% 4. Reference 格式請使用 \href 加上驗證過的 DOI

% Including essential packages for Chinese typesetting
\usepackage{xeCJK}
\usepackage{fontspec}
% Configuring Chinese font with Noto Serif CJK TC, fallback to SimSun
\IfFontExistsTF{Noto Serif CJK TC}{
  \setCJKmainfont{Noto Serif CJK TC}
  \setCJKsansfont{Noto Serif CJK TC}
  \setCJKmonofont{Noto Serif CJK TC}
}{\setCJKmainfont{SimSun}
  \setCJKsansfont{SimSun}
  \setCJKmonofont{SimSun}
}

% Including other necessary packages
\usepackage{geometry}
\usepackage{titlesec}
\usepackage{enumitem}
\usepackage{xcolor}
\usepackage{colortbl}
\usepackage{tcolorbox}
\usepackage{booktabs}
\usepackage{longtable}
\usepackage{array}
\usepackage{graphicx}
\usepackage{tikz}
\usepackage{fancyhdr}
\usepackage{multicol}
\usepackage{datetime2}
\usepackage{hyperref}
\usepackage{mdframed}
\usepackage{amsmath}
\usepackage{amssymb}
\usepackage{multirow}

\hypersetup{
    colorlinks=true,
    linkcolor=emergency_blue,
    citecolor=emergency_blue,
    urlcolor=emergency_blue,
    pdfborder={0 0 0}
}

% Defining page geometry
\geometry{
  top=2.5cm,
  bottom=2.5cm,
  left=2cm,
  right=2cm,
  headheight=25pt
}

% Adjusting topmargin to compensate for increased headheight
\addtolength{\topmargin}{-10pt}

% Defining colors
\definecolor{emergency_red}{RGB}{186,24,27}
\definecolor{emergency_blue}{RGB}{0,114,188}
\definecolor{emergency_gray}{RGB}{120,120,120}
\definecolor{highlight_yellow}{RGB}{255,250,205}

% Setting title formats
\titleformat{\section}[block]{\Large\bfseries\color{emergency_red}}{\thesection}{10pt}{}
\titleformat{\subsection}[block]{\large\bfseries\color{emergency_blue}}{\thesubsection}{8pt}{}
\titleformat{\subsubsection}[block]{\normalsize\bfseries\color{emergency_gray}}{\thesubsubsection}{6pt}{}

% Defining custom environment for important notes
\newmdenv[
    linecolor=emergency_red,
    backgroundcolor=highlight_yellow,
    linewidth=2pt,
    topline=true,
    bottomline=true,
    leftline=true,
    rightline=true,
    innertopmargin=10pt,
    innerbottommargin=10pt,
    innerleftmargin=10pt,
    innerrightmargin=10pt
]{importantbox}

% Defining note command with tcolorbox
\newcommand{\note}[1]{
    \par
    \noindent
    \begin{tcolorbox}[
        colback=emergency_blue!5!white,
        colframe=emergency_blue,
        title=\textbf{重點提醒},
        fonttitle=\bfseries,
        boxrule=0.5pt,
        arc=3pt,
        left=6pt,
        right=6pt,
        top=6pt,
        bottom=6pt
    ]
    #1
    \end{tcolorbox}
}

% Key findings box
\newcommand{\keyfinding}[1]{
    \par
    \noindent
    \begin{tcolorbox}[
        colback=yellow!10!white,
        colframe=orange!75!black,
        title=\textbf{核心發現摘要},
        fonttitle=\bfseries,
        boxrule=0.5pt,
        arc=3pt
    ]
    #1
    \end{tcolorbox}
}

% Defining custom date format for ISO 8601
\DTMnewdatestyle{iso}{
  \renewcommand{\DTMdisplaydate}[4]{\number##1-\number##2-\number##3}
  \renewcommand{\DTMDisplaydate}[4]{\number##1-\number##2-\number##3}
}
\DTMsetdatestyle{iso}
\newcommand{\mydate}{\DTMtoday}


% Custom header for this document
\pagestyle{fancy}
\fancyhf{}
\fancyhead[L]{\textcolor{emergency_red}{心臟磁振造影教學}}
\fancyhead[R]{\textcolor{emergency_blue}{Cardiac MRI - LGE Patterns}}
\fancyfoot[C]{\textcolor{emergency_gray}{\thepage}}
\renewcommand{\headrulewidth}{1pt}
\renewcommand{\headrule}{\hbox to\headwidth{\color{emergency_red}\leaders\hrule height \headrulewidth\hfill}}

\begin{document}

% Title Page
\begin{center}
{\LARGE\bfseries\textcolor{emergency_red}{心臟磁振造影 LGE 判讀教學}}\\[0.3cm]
{\Large\textbf{Late Gadolinium Enhancement Pattern Interpretation}}\\[0.5cm]
{\large Cardiac MRI 心肌病變鑑別診斷指南}\\[0.5cm]
{\normalsize 圖片來源：www.radiologyillustration.com}\\[0.3cm]
{\normalsize 整理：謝慕揚 MD, PhD, FESC \quad 日期：\mydate}
\end{center}

\vspace{0.5cm}

\begin{importantbox}
\textbf{LGE 判讀核心原則}
\begin{itemize}[leftmargin=*]
    \item \textbf{缺血性 (Ischemic)}：從心內膜開始，沿冠狀動脈支配區域分布
    \item \textbf{非缺血性 (Non-ischemic)}：不沿冠狀動脈分布，可為中層、心外膜或瀰漫性
\end{itemize}
\end{importantbox}

%============================================================
\section{Late Gadolinium Enhancement 基礎原理}
%============================================================

\subsection{什麼是 LGE？}

Late Gadolinium Enhancement（延遲顯影強化）是心臟磁振造影中最重要的組織特徵化技術之一。

\begin{longtable}{p{4cm}p{10cm}}
\toprule
\textbf{項目} & \textbf{說明} \\
\midrule
\endfirsthead
\bottomrule
\endfoot
原理\newline (Mechanism) & Gadolinium 對比劑會積聚在細胞外間質擴大的區域，如纖維化、壞死或水腫組織 \\
\midrule
掃描時機\newline (Timing) & 注射對比劑後 10-15 分鐘進行掃描，此時正常心肌已排除對比劑，而病變區域仍有滯留 \\
\midrule
影像表現\newline (Appearance) & 病變區域呈現\textbf{高信號（亮白色）}，正常心肌為\textbf{低信號（暗黑色）} \\
\midrule
臨床意義\newline (Clinical significance) & 偵測心肌纖維化、梗塞、發炎，評估心肌存活性 (viability)，預測預後 \\
\bottomrule
\end{longtable}

\subsection{LGE 判讀流程}

\begin{enumerate}[leftmargin=*]
    \item 確認影像品質：nulling 是否適當，正常心肌應為暗黑色
    \item 辨識強化區域：高信號區域的位置與分布
    \item 判斷強化模式：內膜下、透壁、中層、心外膜、瀰漫性
    \item 對照冠狀動脈分布：是否符合特定冠狀動脈支配區域
    \item 結合臨床病史：缺血性 vs 非缺血性心肌病變
\end{enumerate}

%============================================================
\section{LGE 強化模式總覽}
%============================================================

\begin{center}
\includegraphics[width=0.95\textwidth]{2026-01-11/CMR_LGE_Patterns.png}
\end{center}

\vspace{0.5cm}

%============================================================
\section{缺血性心肌病變 (Ischemic Patterns)}
%============================================================

\note{
\textbf{關鍵特徵}：LGE 從心內膜 (endocardium) 開始，沿冠狀動脈支配區域 (coronary territory) 分布，符合「wavefront phenomenon」（波前現象）
}

\subsection{內膜下梗塞 (Sub-Endocardial Infarct)}

\begin{longtable}{p{4cm}p{10cm}}
\toprule
\textbf{特徵} & \textbf{說明} \\
\midrule
\endfirsthead
\bottomrule
\endfoot
LGE 分布 & 僅侷限於心內膜下區域，未達透壁 \\
\midrule
透壁程度\newline (Transmurality) & $<$50\% 壁厚 \\
\midrule
臨床意義 & 較早期或較小範圍的梗塞\newline 該區域心肌仍有\textbf{存活性 (viability)}，血管重建可能有益 \\
\midrule
鑑別重點 & 必須符合單一冠狀動脈支配區域 \\
\bottomrule
\end{longtable}

\subsection{透壁性梗塞 (Transmural Infarct)}

\begin{longtable}{p{4cm}p{10cm}}
\toprule
\textbf{特徵} & \textbf{說明} \\
\midrule
\endfirsthead
\bottomrule
\endfoot
LGE 分布 & 從心內膜延伸至心外膜，完全穿透心室壁 \\
\midrule
透壁程度\newline (Transmurality) & $>$50\% 壁厚，通常接近 100\% \\
\midrule
臨床意義 & 較嚴重或較晚期的梗塞\newline 該區域心肌\textbf{無存活性}，血管重建無法改善功能 \\
\midrule
預後意義 & LGE 範圍越大，心室功能恢復機會越低，預後越差 \\
\bottomrule
\end{longtable}

\keyfinding{
\textbf{存活性評估 (Viability Assessment)}\\
\begin{itemize}[leftmargin=*]
    \item LGE $<$25\% 壁厚：血管重建後功能恢復機會高
    \item LGE 25-50\% 壁厚：可能部分恢復
    \item LGE $>$50\% 壁厚：功能恢復機會低
\end{itemize}
}

%============================================================
\section{非缺血性心肌病變 (Non-Ischemic Patterns)}
%============================================================

\note{
\textbf{關鍵特徵}：LGE \textbf{不}沿冠狀動脈支配區域分布，可能出現在中層、心外膜、RV 接合點，或呈瀰漫性分布
}

\subsection{中層強化 (Mid-Wall Enhancement)}

\begin{longtable}{p{4cm}p{10cm}}
\toprule
\textbf{相關疾病} & \textbf{特徵說明} \\
\midrule
\endfirsthead
\bottomrule
\endfoot
擴張型心肌病變\newline (Dilated Cardiomyopathy, DCM) & 中隔中層線性強化，常見於 interventricular septum\newline 預後不良指標，與心律不整風險相關 \\
\midrule
心肌炎\newline (Myocarditis) & 中層或心外膜強化\newline 常見於 inferolateral wall\newline 急性期可合併心肌水腫 (T2 高信號) \\
\midrule
類肉瘤\newline (Sarcoidosis) & 多發性、斑塊狀中層或心外膜強化\newline 常伴隨局部室壁運動異常\newline 可能累及右心室 \\
\midrule
Anderson-Fabry 病 & 典型位於 basal inferolateral wall 中層\newline 常合併左心室肥厚 \\
\midrule
Chagas 病 & 心尖部及下壁中層強化\newline 常見於中南美洲流行區 \\
\bottomrule
\end{longtable}

\subsection{心外膜強化 (Epicardial Enhancement)}

\begin{longtable}{p{4cm}p{10cm}}
\toprule
\textbf{相關疾病} & \textbf{特徵說明} \\
\midrule
\endfirsthead
\bottomrule
\endfoot
心肌炎\newline (Myocarditis) & 心外膜或中層強化\newline 典型位於 inferolateral wall\newline 病毒性心肌炎最常見 \\
\midrule
類肉瘤\newline (Sarcoidosis) & 可呈心外膜或中層分布\newline 常為多發性 \\
\midrule
Anderson-Fabry 病 & 可延伸至心外膜 \\
\midrule
Chagas 病 & 心外膜或中層分布 \\
\bottomrule
\end{longtable}

\subsection{右心室接合點強化 (RV Insertion Point Enhancement)}

\begin{longtable}{p{4cm}p{10cm}}
\toprule
\textbf{相關疾病} & \textbf{特徵說明} \\
\midrule
\endfirsthead
\bottomrule
\endfoot
肥厚型心肌病變\newline (HCM) & 前後 RV insertion points 強化\newline 與纖維化程度相關\newline 可能與心律不整風險有關 \\
\midrule
右心室壓力過負荷\newline (RV Pressure Overload) & 常見於 Pulmonary hypertension\newline 反映 RV-LV 交界處機械應力增加 \\
\bottomrule
\end{longtable}

\subsection{瀰漫性/全域性強化 (Diffuse/Global Enhancement)}

\begin{longtable}{p{4cm}p{10cm}}
\toprule
\textbf{相關疾病} & \textbf{特徵說明} \\
\midrule
\endfirsthead
\bottomrule
\endfoot
心臟類澱粉沉積症\newline (Cardiac Amyloidosis) & 瀰漫性心內膜下或透壁強化\newline 傳統 LGE 技術可能難以偵測（因正常心肌也受影響）\newline \textbf{T1 mapping} 更為敏感：native T1 顯著升高 \\
\midrule
全身性硬化症\newline (Systemic Sclerosis) & 瀰漫性線性或斑塊狀強化\newline 常合併心包積液 \\
\midrule
心臟移植後\newline (Post-transplant) & 可能出現瀰漫性強化\newline 反映慢性排斥或心肌病變 \\
\bottomrule
\end{longtable}

%============================================================
\section{LGE 模式鑑別診斷總表}
%============================================================

\begin{longtable}{p{3.5cm}p{4cm}p{6cm}}
\toprule
\textbf{強化模式} & \textbf{分布特徵} & \textbf{常見疾病} \\
\midrule
\endfirsthead
\toprule
\textbf{強化模式} & \textbf{分布特徵} & \textbf{常見疾病} \\
\midrule
\endhead
\bottomrule
\endfoot
\multicolumn{3}{c}{\textbf{缺血性 (Ischemic)}} \\
\midrule
內膜下\newline (Subendocardial) & 沿冠狀動脈分布\newline $<$50\% 壁厚 & 急性/亞急性 MI\newline 有存活性 \\
\midrule
透壁性\newline (Transmural) & 沿冠狀動脈分布\newline $>$50\% 壁厚 & 陳舊性 MI\newline 無存活性 \\
\midrule
\multicolumn{3}{c}{\textbf{非缺血性 (Non-Ischemic)}} \\
\midrule
中層\newline (Mid-wall) & Septum 或其他位置\newline 線性分布 & DCM, Myocarditis,\newline Sarcoidosis, Anderson-Fabry \\
\midrule
心外膜\newline (Epicardial) & Inferolateral wall 常見 & Myocarditis, Sarcoidosis \\
\midrule
RV 接合點\newline (RV insertion) & 前後 RV-LV 接合處 & HCM, Pulmonary HTN \\
\midrule
瀰漫性\newline (Diffuse) & 全心室，難以 null & Amyloidosis,\newline Systemic sclerosis \\
\bottomrule
\end{longtable}

%============================================================
\section{重點整理與臨床建議}
%============================================================

\begin{importantbox}
\textbf{LGE 判讀重點}
\begin{enumerate}[leftmargin=*]
    \item \textbf{先判斷缺血 vs 非缺血}：內膜下起始 + 冠狀動脈分布 = 缺血性
    \item \textbf{評估透壁程度}：影響存活性評估與治療決策
    \item \textbf{結合臨床病史}：胸痛史、心肌酵素、冠狀動脈造影結果
    \item \textbf{注意特殊模式}：中層強化考慮 DCM/myocarditis，RV insertion 考慮 HCM
    \item \textbf{搭配其他序列}：T1/T2 mapping, STIR 等進一步鑑別
\end{enumerate}
\end{importantbox}

\subsection{何時懷疑非缺血性病因？}

\begin{itemize}[leftmargin=*]
    \item LGE 不沿冠狀動脈分布
    \item 年輕患者無明顯心血管危險因子
    \item 強化位於中層或心外膜
    \item 多發性斑塊狀強化
    \item 冠狀動脈造影正常
\end{itemize}

\subsection{進階影像技術}

\begin{longtable}{p{4cm}p{10cm}}
\toprule
\textbf{技術} & \textbf{臨床應用} \\
\midrule
\endfirsthead
\bottomrule
\endfoot
T1 Mapping & 偵測瀰漫性纖維化，對 amyloidosis 特別敏感\newline Native T1 升高提示間質擴張 \\
\midrule
T2 Mapping / STIR & 偵測心肌水腫，區分急性 vs 慢性病變\newline 急性心肌炎診斷重要依據 \\
\midrule
ECV Mapping & 量化細胞外容積比例\newline 更精確評估瀰漫性纖維化 \\
\bottomrule
\end{longtable}

\vfill
\begin{center}
\textcolor{emergency_gray}{\small 本教材僅供教學參考，臨床判讀請結合完整檢查與病史}\\
\textcolor{emergency_gray}{\small 原始圖片來源：www.radiologyillustration.com}
\end{center}

\end{document}
